\documentclass[10pt,a4paper]{amsart}
\usepackage[margin=19mm]{geometry}
\usepackage{amsmath,amssymb,mathtools}
\usepackage[scr=rsfs,cal=euler]{mathalfa}
\usepackage{xfrac}
\usepackage[unicode,hidelinks,bookmarks=false]{hyperref}
\newcommand{\ie}{i.e.\ }
\newcommand{\cf}{cf.\ }
\newcommand{\resp}{respectively\ }
\newcommand{\Subsection}{\subsection}
\providecommand{\nobreakdash}{\nobreak\hbox{-}}
\setlength{\emergencystretch}{2em}
\multlinegap0pt
\usepackage{bm}
\usepackage{bbm}
\usepackage{dsfont}
\usepackage{xspace}
\usepackage{cancel}
\usepackage{mathtools}
\usepackage[shortlabels]{enumitem}
\usepackage{amsthm}
\makeatletter
\@ifundefined{theorem}{\newtheorem{theorem}{Theorem}}{}
\@ifundefined{algorithm}{\newtheorem{algorithm}[theorem]{Algorithm}}{}
\@ifundefined{definition}{\newtheorem{definition}[theorem]{Definition}}{}
\@ifundefined{lemma}{\newtheorem{lemma}[theorem]{Lemma}}{}
\@ifundefined{proposition}{\newtheorem{proposition}[theorem]{Proposition}}{}
\makeatother
\newcommand{\Gc}{\mathcal G}
\newcommand{\Lc}{\mathcal L}
\newcommand{\Prob}{\mathbb P}
\newcommand{\supp}{\operatorname{supp}}
\title[The Brownian loop-catcher]{The Brownian loop-catcher}
\author{Gefei Cai}
\date{Source: arXiv:2607.18070v2 (CC BY 4.0)}
\begin{document}
\maketitle

\noindent\emph{Source and licence.} The Brownian loop-catcher. Version arXiv:2607.18070v2 (CC BY 4.0), distributed under the Creative Commons Attribution 4.0 International licence, as recorded in the arXiv licence metadata for this exact version and shown on the abstract page https://arxiv.org/abs/2607.18070v2. Full rights notice: licenses/arxiv-2607.18070-CC-BY-4.0.txt.\par\smallskip

\noindent\textit{Editorial scope.} Counted unit: the Proposition whose source label is \texttt{prop:path-recovery} in the article (source0.tex lines 1066--1072, source label \texttt{prop:path-recovery}), delivered together with the article's own complete original proof of it (source0.tex lines 1073--1157, one \texttt{proof} environment of 85 lines). No mathematics is written, repaired or re-derived: statement and proof are the article's own text line for line. Required context retained uncounted: def:entangled (lines 776--785); eq:sequential-reduction (lines 778--782); prop:entangled-recovery (lines 791--799); lem:one-step (lines 804--835); eq:rooted-walk-law (lines 831--834); alg:time-catcher (lines 944--993); alg:1 (lines 947--989); alg:2 (lines 958--963); alg:3 (lines 958--963); eq:initial-block-law (lines 958--963); alg:time-recovery (lines 1038--1064). Every definition, display and label of those lines is retained unchanged. Omitted from this package and not redistributed: 1--775, 783--790, 800--803, 835--943, 964--1037, 1065--1065, 1158--3556. Nothing inside a retained excerpt is omitted or altered: every retained body is delivered exactly as the article's lines are written, so no omission receipt of any kind accompanies this package. Citations: the retained excerpts carry no citation command at all, so no bibliography entry of the article is transcribed and no reference is redistributed. Reference adaptation: none was needed and none was made; every reference inside the delivered unit resolves inside it, so no unresolved reference or citation is produced. Printed slips kept literal: no printed word, symbol, hypothesis or number of the counted statement or of its proof was altered. Layout and class: the article's own class is \texttt{article} with packages including amsmath, amssymb, amscd, amsthm, amsfonts, bbm, mathrsfs, mathtools, booktabs, longtable, array, microtype, graphicx, float; the wrapper is the lane's pinned \texttt{amsart} editorial wrapper at 10pt on A4, which restates the theorem-like environments the retained text needs and adds the article's own macro definitions that the retained text needs (\texttt{\textbackslash Gc}, \texttt{\textbackslash Lc}, \texttt{\textbackslash Prob}, \texttt{\textbackslash supp}), with extra wrapper packages (bm, bbm, dsfont, xspace, cancel, mathtools, enumitem). The delivered numbering is the unit's own. Assets: no retained span contains a figure environment, a graphics-inclusion command, a PDF-page inclusion, a table or a TikZ picture; no image file is copied into this package, the package declares no asset dependency and no image file is redistributed with it. Licence: version arXiv:2607.18070v2 of this article is distributed under the Creative Commons Attribution 4.0 International licence, as recorded in the arXiv licence metadata for this exact version and shown on the abstract page https://arxiv.org/abs/2607.18070v2. A full rights notice is delivered at licenses/arxiv-2607.18070-CC-BY-4.0.txt. Characters: every retained excerpt is pure ASCII, so the delivered engine, XeLaTeX with its default Latin Modern text font, reports no missing character.
\begin{definition}[entangled multipath loop-erasure]\label{def:entangled}
For $r\ge1$ and an ordered tuple of directed paths $\omega_i:a_i\to b_i$, $1\le i\le r$, put
\begin{equation}\label{eq:sequential-reduction}
 H_i=P(\omega_i),\qquad C_0=\varnothing,\qquad
 P_i=\Phi_{C_{i-1}}(H_i),\qquad
 C_i=C_{i-1}\cup\supp(P_i).
\end{equation}
Let $\gamma_i$ be the directed path corresponding to the record $P_i$ (equipped with the time parametrization).  We call $(\gamma_1,\ldots,\gamma_r)$
the entangled multipath loop-erasure of $(\omega_1,\ldots,\omega_r)$, whose range is $C_r$.
\end{definition}

\begin{proposition}\label{prop:entangled-recovery}
Let $(P_1,\ldots,P_r)$ be the output records of independent normalized random walks
$\omega_i:a_i\to b_i$ ($1\le i\le r$) in Definition~\ref{def:entangled}, and let $\gamma_i$ be the corresponding
 entangled multipath LERW.  Independently sample an unrooted random-walk loop soup in $R$ with intensity $1$, and assign each loop that intersects $C_r$ to the smallest $i$ for which
it intersects $\gamma_i$, and then to the earliest path time $s$ on $\gamma_i$
at which such an intersection occurs.  If the loop visits $\gamma_i(s)$ more than once, choose uniformly one of those visits as its starting time. For each time on the entangled multipath LERW, order its assigned random-walk loops by an independent
uniform permutation and insert them in that order. Then conditioned on
$(P_1,\ldots,P_r)$, the resulting tuple of paths has the same law as the conditional law of $(\omega_1,\ldots,\omega_r)$ given $(P_1,\ldots,P_r)$.
\end{proposition}

\begin{lemma}\label{lem:one-step}
Fix $A\subset R$,  let $H$ be a legal record from $a$ to $b$, and set
$P=\Phi_A(H)$.  List $\supp(P)\setminus A$ as $u_1,\ldots,u_m$ in first-visit
order along the directed path $\gamma$ corresponding to $P$. For each $1\le j\le m$, let $t_j$ be the first-visit time of $u_j$
and
\[
 D_j=R\setminus\bigl(A\cup\{u_1,\ldots,u_{j-1}\}\bigr).
\]
Then there are unique ordered lists $(\ell_{j,s})_{1\le s\le n_j}$ of loops in
$D_j$ from $u_j$ to $u_j$, with no intermediate visit to $u_j$, whose insertion
in $\gamma$ at $t_j$ reconstructs $H$.  Conversely, every such collection of lists determines one legal
record $H$ with $\Phi_A(H)=P$, and
\begin{equation}\label{eq:positive-fiber-weight}
 q(H)=q(P)\prod_{j=1}^m\prod_{s=1}^{n_j}q(\ell_{j,s}).
\end{equation}
Consequently,
\begin{align}
 d_{R\setminus A}
 \sum_{\substack{H\text{ legal in }R\\\Phi_A(H)=P}}q(H)
 &=q(P)d_{R\setminus(A\cup\supp(P))},\label{eq:fiber}\\
 \Gc_R(a,b)d_{R\setminus A}
 &=\sum_{\substack{P:a\to b\text{ legal in }R\\\Phi_A(P)=P}}
 q(P)d_{R\setminus(A\cup\supp(P))}.\label{eq:one-step}
\end{align}

As a corollary, suppose $H$ is the record of a normalized random walk from $a$ to $b$ in $R$. Then conditioned on $P=\Phi_A(H)$, the lists are
independent and
\begin{equation}\label{eq:rooted-walk-law}
 \Prob[(\ell_{j,1},\ldots,\ell_{j,n_j})=(\xi_1,\ldots,\xi_d)\mid P]
 =\frac{\prod_{s=1}^d q(\xi_s)}{\Gc_{D_j}(u_j,u_j)}.
\end{equation}
\end{lemma}

\begin{algorithm}[Construction]
\label{alg:time-catcher} Let $(\omega_1,\ldots,\omega_r)$ be an ordered tuple of directed paths in $R\subset V$.

\begin{enumerate}[label=\textup{(\arabic*)},leftmargin=2.4em]
\item Let $(P_i,C_i)$ be as in~\eqref{eq:sequential-reduction}; let $\gamma_i$
be the path corresponding to $P_i$, and put
\[
 S=C_r,\qquad B=R\setminus S.
\]
These $\gamma_i$ form the current layer.\label{alg:1}

\item At each first-visit time in Lemma~\ref{lem:one-step}, let
$\ell_1,\ldots,\ell_d$ be the uniquely determined ordered list of loops.  Independently
at different first-visit times, choose $K\in\{0,\ldots,d\}$ with
\begin{equation}\label{eq:initial-block-law}
 \Prob[K=k]
 =\frac{(\lambda)^{\overline k}}{k!}
  \frac{(1-\lambda)^{\overline{d-k}}}{(d-k)!},
 \qquad 0\le k\le d,
\end{equation}
which is the beta--binomial distribution with parameters $d$, $\lambda$, and $1-\lambda$.
Delete $\ell_1,\ldots,\ell_K$ and retain
$\ell_{K+1},\ldots,\ell_d$.\label{alg:2}

\item For every retained loop $\ell$, list in its time order all maximal
segments
\[
 x,b_0,\ldots,b_n,y,
 \qquad x,y\in S,\quad b_0,\ldots,b_n\in B,\quad n\ge0.
\]
View $(b_0,\ldots,b_n)$ as a directed path in $B$. Label
this path by $(i,t,h,s)$, where $i$ is the index of the current path
$\gamma_i$, $t$ is the first-visit time of $\gamma_i$ at which $\ell$ is
attached, $h$ is the index of $\ell$ in the ordered list of loops at time $t$, and $s$ is the starting time of the segment within $\ell$. Order the paths $(b_0,\ldots,b_n)$ lexicographically by these
labels. Denote the resulting ordered tuple by
$(\zeta_1,\ldots,\zeta_N)$.\footnote{The portions of each retained loop lying in $S$ are left unchanged. In particular, a retained loop contained entirely in $S$ contributes no path to $(\zeta_1,\ldots,\zeta_N)$.}\label{alg:3}

\item If $N>0$, apply this algorithm recursively to the ordered
tuple $(\zeta_1,\ldots,\zeta_N)$ in $B$.  Denote the returned paths by
$(\zeta_1^\lambda,\ldots,\zeta_N^\lambda)$.  Replace each $\zeta_j$ by
$\zeta_j^\lambda$ at its recorded position.  If $N=0$, make no recursive
call. Note that the recursion terminates since $B$ is a proper subset of $R$.

\item After these replacements, insert the resulting retained loops (in their recorded order) at the corresponding first-visit times of the
$\gamma_i$.
\end{enumerate}

The output is a tuple of directed paths
$\omega_1^\lambda,\ldots,\omega_r^\lambda$. We also set $S_{\lambda,\times}:=\bigcup_{i=1}^r\supp(\omega_i^\lambda)$.
\end{algorithm}

\begin{algorithm}[Recovery]
\label{alg:time-recovery}
Let $\Omega^{(0)}=(\omega_1^\lambda,\ldots,\omega_r^\lambda)$ be the output of
Algorithm~\ref{alg:time-catcher}, and let
$\Lc_R^\lambda$ be a realization of the random-walk loop soup in $R$ with intensity $\lambda$.

\begin{enumerate}[label=\textup{(\arabic*)},leftmargin=2.4em]
\item Recover all recursive layers $0\le n\le L$ as above.

\item For each $\ell\in\Lc_R^\lambda$ that intersects
$S_{\lambda,\times}$, let $n$ be the smallest index such that
$\ell\cap S_n\ne\varnothing$. Within layer~$n$, choose the smallest path index $i$ for
which $\ell$ meets $\gamma_i^{(n)}$, and then the earliest path time $t$ at
which it is met. If $\ell$ visits $\gamma_i^{(n)}(t)$ more than once, root
it uniformly at one of those visits. At each fixed $(n,i,t)$, order the
assigned rooted loops by an independent uniform permutation.

\item Process the layers in the order $L,L-1,\ldots,0$. At each
first-visit time, prepend the assigned rooted loops (understood as cut at their successive returns to the assigned root), in their chosen order,
to the ordered list of loops already there, and insert the resulting list as
in Lemma~\ref{lem:one-step}. After completing layer~$n>0$, substitute the
resulting paths into their recorded positions in the loops retained at
layer~$n-1$.
\end{enumerate}

The output is the directed paths obtained at layer~$0$.
\end{algorithm}

\begin{proposition}\label{prop:path-recovery}
Fix $R\subset V$, ordered endpoint pairs $(a_i,b_i)_{1\le i\le r}$ with
$\Gc_R(a_i,b_i)>0$, and $0\le\lambda\le 1$. Let
$\omega_i:a_i\to b_i$ be independent normalized random walks in $R$, and let $\Omega^{(0)}
 =(\omega_1^\lambda,\ldots,\omega_r^\lambda)$ be the output of Algorithm~\ref{alg:time-catcher}. Sample an independent random-walk loop soup $\Lc_R^\lambda$ in $R$ with intensity $\lambda$. Then the output of Algorithm~\ref{alg:time-recovery} has
the same joint law as $(\omega_1,\ldots,\omega_r)$.
\end{proposition}

\begin{proof}
We first record a loop-soup calculation.  Fix $u\in D\subset R$ and write $G=\Gc_D(u,u)$.
Consider an intensity-$\alpha$ random-walk loop soup in $D$.  Root each
loop visiting $u$ uniformly at one of its visits to $u$, uniformly order
the rooted loops, and cut them at their successive returns to $u$.  For
an ordered list $(\xi_1,\ldots,\xi_d)$ of loops that have multiplicity $1$ at $u$, the probability of obtaining this list is
\begin{equation}\label{eq:alpha-root-list}
 \Prob\bigl[(\xi_1,\ldots,\xi_d)\bigr]
 =
 G^{-\alpha}\frac{(\alpha)^{\overline d}}{d!}
 \prod_{s=1}^d q(\xi_s).
\end{equation}
Indeed, the loop mass of random-walk loops in $D$ visiting $u$ is $\log G$, and if
the list is produced by $M$ random-walk loops and the $m$-th loop is cut into
$j_m$ consecutive paths, then $j_1+\cdots+j_M=d$
and the corresponding contribution is $ G^{-\alpha}
 \frac{\alpha^M}{M!j_1\cdots j_M}
 \prod_{s=1}^d q(\xi_s)$.
Summing over $M$ and the ordered compositions of $d$ gives~\eqref{eq:alpha-root-list} (doing the same calculation as in the proof of Proposition~\ref{prop:entangled-recovery}).

We next apply~\eqref{eq:alpha-root-list} to Step~\ref{alg:2} of
Algorithm~\ref{alg:time-catcher}.  For an ordered tuple of independent normalized random walks, apply Algorithm~\ref{alg:time-catcher} and let $\boldsymbol P=(P_1,\ldots,P_m)$
be the records obtained in Step~\ref{alg:1}.  Conditioned on
$\boldsymbol P$, Lemma~\ref{lem:one-step}, applied successively in path
order, gives independent ordered lists of loops at the first-visit times.
Fix a first-visit time, let $u$ and $D$ be its vertex and associated domain
in Lemma~\ref{lem:one-step}, and let $(\ell_1,\ldots,\ell_d)$
be the corresponding ordered list of loops at the beginning of Step~\ref{alg:2}.  For
$0\le k\le d$, \eqref{eq:rooted-walk-law} and
\eqref{eq:initial-block-law} imply that the conditional probability $\Prob\!\left[
   (\ell_1,\ldots,\ell_d)=(\xi_1,\ldots,\xi_d),\,K=k
   \,\middle|\,\boldsymbol P
 \right]$ is equal to
\begin{align}\label{eq:list-splitting}
 \frac{\prod_{s=1}^d q(\xi_s)}{G}
 \frac{(\lambda)^{\overline k}}{k!}
 \frac{(1-\lambda)^{\overline{d-k}}}{(d-k)!}=
 \left[
  G^{-\lambda}
  \frac{(\lambda)^{\overline k}}{k!}
  \prod_{s=1}^k q(\xi_s)
 \right]
 \left[
  G^{-(1-\lambda)}
  \frac{(1-\lambda)^{\overline{d-k}}}{(d-k)!}
  \prod_{s=k+1}^d q(\xi_s)
 \right].
\end{align}
By \eqref{eq:alpha-root-list}, the two factors on the right side of~\eqref{eq:list-splitting} are the probabilities of obtaining $(\xi_1,\ldots,\xi_k)$ and $(\xi_{k+1},\ldots,\xi_d)$
from independent loop soups of intensities $\lambda$ and $1-\lambda$,
respectively.  Hence, conditioned on $\boldsymbol P$, the list deleted in
Step~\ref{alg:2} is independent of the retained list and has the law
obtained by an independent intensity-$\lambda$ loop soup.  This holds
jointly at all first-visit times.

We now prove the proposition by induction on $|R|$.  Let
\[
 S=S_0,\qquad B=R\setminus S,
\]
and let $(\zeta_1,\ldots,\zeta_N)$ be the ordered tuple constructed in
Step~\ref{alg:3} of Algorithm~\ref{alg:time-catcher}.  Conditioned on the
portions of the retained loops outside these maximal subpaths, their
conditional weight is proportional to $\prod_{j=1}^N q(\zeta_j)$.
Consequently, conditioned on their endpoints, the $(\zeta_j)_{1\le j\le N}$ are independent
normalized random walks in $B$.
The restriction of $\Lc_R^\lambda$ to loops contained in $B$ is an
intensity-$\lambda$ random-walk loop soup in $B$. If $N>0$, then
$B\subsetneq R$, so the induction hypothesis applied in $B$ shows that Algorithm~\ref{alg:time-recovery} recovers the joint
law of $(\zeta_1,\ldots,\zeta_N)$.  Replacing the recovered paths at their
recorded positions therefore recovers the retained loops at layer~$0$.


It remains to recover the lists deleted at recursive layer~$0$.  The loops of
$\Lc_R^\lambda$ intersecting $S$ are assigned by
Algorithm~\ref{alg:time-recovery} according to the same rule as in
Proposition~\ref{prop:entangled-recovery}.  By
\eqref{eq:list-splitting}, its ordered list on each first-visit
time has the conditional law of the list deleted in
Step~\ref{alg:2}, independently of the retained list. Prepending it therefore recovers the full list in
Lemma~\ref{lem:one-step}.
By Proposition~\ref{prop:entangled-recovery}, conditioned on
$\boldsymbol P$, the paths obtained at layer~$0$ have the conditional
joint law of $(\omega_1,\ldots,\omega_r)$.  Averaging over
$\boldsymbol P$ then yields the result.
\end{proof}

\end{document}
